\documentclass[12pt]{article}
\usepackage{amsmath,amssymb,amsfonts,amsthm}
\usepackage[margin=1in]{geometry}

\begin{document}

\begin{center}
{\Large\bfseries Chapter 4, Problem 3}\\[6pt]
Byron \& Fuller, \textit{Mathematics of Classical and Quantum Physics}
\end{center}

\bigskip

\textbf{Problem.} Prove that the $n \times n$ matrix $(I - iA)$ has an inverse if $A$ is Hermitian.

\bigskip
\textbf{Solution.}

A matrix is invertible if and only if its determinant is nonzero, which is equivalent to saying it has no zero eigenvalue. We will show that $(I - iA)$ has no zero eigenvalue.

Suppose $(I - iA)x = 0$ for some vector $x$. Then $x = iAx$, i.e., $Ax = -ix$, so $x$ is an eigenvector of $A$ with eigenvalue $-i$.

But since $A$ is Hermitian, all eigenvalues of $A$ are real. Therefore $-i$ cannot be an eigenvalue of $A$, which means no nonzero $x$ can satisfy $(I - iA)x = 0$.

Hence $(I - iA)$ is nonsingular and therefore invertible. \qed

\bigskip
\textbf{Alternative proof using inner products.}

Suppose $(I - iA)x = 0$. Take the inner product with $x$:
\[
(x, (I - iA)x) = (x, x) - i(x, Ax) = 0.
\]
Since $A$ is Hermitian, $(x, Ax)$ is real. Also $(x, x) = \|x\|^2$ is real and non-negative. Separating into real and imaginary parts:
\[
\|x\|^2 = 0 \quad \text{and} \quad (x, Ax) = 0.
\]
From $\|x\|^2 = 0$ we get $x = 0$. Therefore $(I - iA)$ has trivial kernel and is invertible. \qed

\end{document}
